The gas-generation subsystem is coupled to internal pressure and venting
models in order to represent cell pressurization and subsequent gas
discharge.

The current formulation combines an ideal-gas pressure model with
compressible single-phase vent flow.


\subsection{Ideal-gas pressure}

The cell gas pressure is calculated using

\begin{equation}
    P
    V_{\mathrm{free}}
    =
    n_{\mathrm{tot}}
    R
    T,
\end{equation}

or

\begin{equation}
    P
    =
    \frac{
        n_{\mathrm{tot}}RT
    }{
        V_{\mathrm{free}}
    },
    \label{eq:pressure_ideal}
\end{equation}

where $V_{\mathrm{free}}$ is the available gas volume.

All pressures are treated as absolute pressures.


\subsection{Initial gas inventory}

The initial gas amount is calculated from the specified initial pressure
$P_0$, initial temperature $T_0$, and free volume:

\begin{equation}
    n_0
    =
    \frac{
        P_0 V_{\mathrm{free}}
    }{
        R T_0
    }.
    \label{eq:initial_gas_moles}
\end{equation}

Before venting, the total amount of gas can be represented as the sum of
the initial inventory and generated gas.

During venting, however, pressure must be reconstructed directly from the
instantaneous integrated species inventories.


\subsection{Pressure-temperature coupling}

For fixed volume,

\begin{equation}
    P
    \propto
    n_{\mathrm{tot}} T.
\end{equation}

Cell pressure can therefore rise because of both

\begin{enumerate}
    \item gas generation, which increases $n_{\mathrm{tot}}$;
    \item thermal runaway, which increases $T$.
\end{enumerate}

The pressure response is thus strongly coupled to the thermal and
reaction subsystems.


\subsection{Vent-opening criterion}

A vent-opening pressure

\begin{equation}
    P_{\mathrm{vent}}
\end{equation}

may be specified as part of the thermal problem.

The vent-opening event occurs when

\begin{equation}
    P
    =
    P_{\mathrm{vent}}.
    \label{eq:vent_open_condition}
\end{equation}

Before this condition is reached, the gas system is treated as closed.

After the event, the vent is considered permanently open.


\subsection{Maximum-pressure criterion}

An independent maximum-pressure threshold may also be specified,

\begin{equation}
    P_{\max}.
\end{equation}

This condition may be used to terminate numerical integration when

\begin{equation}
    P
    =
    P_{\max}.
    \label{eq:max_pressure_condition}
\end{equation}

The maximum-pressure threshold is distinct from the vent-opening
pressure and may represent a numerical safety criterion or a simplified
containment-failure limit.


\subsection{Compressible vent flow}

Gas discharge through the vent is represented using compressible
ideal-gas flow.

The upstream pressure is

\begin{equation}
    P_u
    =
    P_{\mathrm{cell}},
\end{equation}

while the downstream pressure is typically the ambient pressure,

\begin{equation}
    P_d
    =
    P_{\mathrm{ambient}}.
\end{equation}

The pressure ratio is

\begin{equation}
    r_p
    =
    \frac{
        P_d
    }{
        P_u
    }.
    \label{eq:pressure_ratio}
\end{equation}


\subsection{Specific gas constant}

The specific gas constant of the mixture is calculated as

\begin{equation}
    R_s
    =
    \frac{
        R
    }{
        \overline{M}
    },
    \label{eq:specific_gas_constant}
\end{equation}

where $\overline{M}$ is the instantaneous mean molar mass of the gas
mixture.


\subsection{Heat-capacity ratio}

The current vent model uses a prescribed heat-capacity ratio

\begin{equation}
    \gamma
    =
    \frac{
        c_p
    }{
        c_v
    }.
\end{equation}

In the present implementation, $\gamma$ is treated as constant rather
than being recomputed from the instantaneous gas composition.


\subsection{Critical pressure ratio}

For ideal-gas isentropic flow, the critical downstream-to-upstream
pressure ratio is

\begin{equation}
    r_{\mathrm{crit}}
    =
    \left(
        \frac{
            2
        }{
            \gamma+1
        }
    \right)^{
        \frac{
            \gamma
        }{
            \gamma-1
        }
    }.
    \label{eq:critical_pressure_ratio}
\end{equation}

The flow is considered choked when

\begin{equation}
    \frac{P_d}{P_u}
    \leq
    r_{\mathrm{crit}}.
    \label{eq:choked_condition}
\end{equation}


\subsection{Choked mass flow}

For choked flow, the mass discharge rate is

\begin{equation}
    \dot{m}_{\mathrm{vent}}
    =
    C_d A_v P_u
    \sqrt{
        \frac{\gamma}{R_s T}
    }
    \left(
        \frac{2}{\gamma+1}
    \right)^{
        \frac{\gamma+1}{2(\gamma-1)}
    }.
    \label{eq:choked_mass_flow}
\end{equation}

where $C_d$ is the discharge coefficient and $A_v$ is the effective vent
area.

Once the flow is choked, decreasing downstream pressure further does not
increase the mass flow predicted by this relation.


\subsection{Unchoked mass flow}

When

\begin{equation}
    \frac{P_d}{P_u}
    >
    r_{\mathrm{crit}},
\end{equation}

the discharge is treated as unchoked.

The mass flow rate is

\begin{equation}
    \dot{m}_{\mathrm{vent}}
    =
    C_d A_v P_u
    \sqrt{
        \frac{2\gamma}
             {R_s T(\gamma-1)}
        \left[
            \left(
                \frac{P_d}{P_u}
            \right)^{2/\gamma}
            -
            \left(
                \frac{P_d}{P_u}
            \right)^{(\gamma+1)/\gamma}
        \right]
    }.
    \label{eq:unchoked_mass_flow}
\end{equation}

As $P_d$ approaches $P_u$, the predicted discharge rate approaches zero.


\subsection{Reverse flow}

The current formulation does not represent inward flow through the vent.

Consequently,

\begin{equation}
    P_u
    \leq
    P_d
    \quad
    \Longrightarrow
    \quad
    \dot{m}_{\mathrm{vent}}
    =
    0.
    \label{eq:no_reverse_flow}
\end{equation}


\subsection{Molar vent flow}

The total mass discharge rate is converted to a molar flow rate using

\begin{equation}
    \dot{n}_{\mathrm{vent}}
    =
    \frac{
        \dot{m}_{\mathrm{vent}}
    }{
        \overline{M}
    }.
    \label{eq:total_molar_vent_flow}
\end{equation}


\subsection{Species vent rates}

The current model assumes a well-mixed gas inventory.

The vented gas therefore has the same mole fractions as the gas
remaining inside the modeled free volume.

For species $s$,

\begin{equation}
    \dot{n}_{s,\mathrm{vent}}
    =
    x_s
    \dot{n}_{\mathrm{vent}}.
    \label{eq:species_vent_flow}
\end{equation}

Summing over species gives

\begin{equation}
    \sum_s
    \dot{n}_{s,\mathrm{vent}}
    =
    \dot{n}_{\mathrm{vent}}.
\end{equation}


\subsection{Closed and open governing equations}

The vent model creates two distinct gas-balance phases.

Before vent opening,

\begin{equation}
    \frac{dn_s}{dt}
    =
    \dot{n}_{s,\mathrm{gen}}.
    \label{eq:vent_closed_phase}
\end{equation}

After vent opening,

\begin{equation}
    \frac{dn_s}{dt}
    =
    \dot{n}_{s,\mathrm{gen}}
    -
    \dot{n}_{s,\mathrm{vent}}.
    \label{eq:vent_open_phase}
\end{equation}


\subsection{Irreversible vent transition}

The current LiionTR implementation treats vent opening as an irreversible
state transition,

\begin{equation}
    \mathrm{CLOSED}
    \longrightarrow
    \mathrm{OPEN}.
\end{equation}

The vent does not reseal if pressure subsequently falls below
$P_{\mathrm{vent}}$.

The numerical interpretation of this transition is illustrated in
\cref{fig:vent_state_machine}.

\begin{figure}[htbp]
    \centering

    \begin{tikzpicture}[
        node distance=4.0cm,
        state/.style={
            rectangle,
            rounded corners,
            draw,
            align=center,
            minimum width=3.5cm,
            minimum height=1.3cm
        },
        line/.style={
            -{Latex[length=3mm]},
            thick
        }
    ]

    \node[state] (closed)
        {Vent closed\\
         $\dot{n}_s=\dot{n}_{s,\mathrm{gen}}$};

    \node[state, right=of closed] (open)
        {Vent open\\
         $\dot{n}_s=
         \dot{n}_{s,\mathrm{gen}}
         -
         \dot{n}_{s,\mathrm{vent}}$};

    \draw[line]
        (closed)
        --
        node[above, align=center]
        {$P=P_{\mathrm{vent}}$}
        (open);

    \end{tikzpicture}

    \caption{
        Irreversible vent-state transition used by the current LiionTR
        solver.
    }
    \label{fig:vent_state_machine}
\end{figure}


\subsection{Coupled pressure-vent feedback}

After vent opening, pressure determines the vent-flow rate, while the
vent-flow rate modifies the gas inventory and therefore pressure.

This feedback may be represented as

\begin{equation}
    P
    \longrightarrow
    \dot{n}_{\mathrm{vent}}
    \longrightarrow
    n_{\mathrm{tot}}
    \longrightarrow
    P.
\end{equation}

Gas generation and temperature continue to evolve simultaneously.

The pressure response therefore depends on the competition between

\begin{equation}
    \dot{n}_{\mathrm{gen}}
\end{equation}

and

\begin{equation}
    \dot{n}_{\mathrm{vent}},
\end{equation}

together with the evolving thermal state.


\subsection{Current assumptions}

The pressure and venting formulation assumes:

\begin{itemize}
    \item ideal-gas behavior;
    \item constant free volume;
    \item uniform gas pressure;
    \item gas temperature equal to the lumped cell temperature;
    \item single-phase gas discharge;
    \item quasi-steady compressible vent flow;
    \item fixed effective vent area;
    \item constant discharge coefficient;
    \item constant heat-capacity ratio;
    \item a well-mixed gas inventory;
    \item no preferential species transport through the vent;
    \item no reverse flow;
    \item irreversible vent opening.
\end{itemize}


\subsection{Limitations}

Real lithium-ion battery venting can involve considerably more complex
behavior than represented by the present model.

The discharged material may contain

\begin{itemize}
    \item permanent gases;
    \item electrolyte vapor;
    \item liquid droplets;
    \item aerosols;
    \item condensed particles;
    \item decomposed electrode and separator material.
\end{itemize}

The physical vent may also deform progressively, resulting in a
time-dependent effective flow area.

These mechanisms are outside the scope of the present reduced-order
model.


\subsection{Future extensions}

Future developments may include:

\begin{itemize}
    \item gas properties obtained dynamically from Cantera;
    \item temperature- and composition-dependent $\gamma$;
    \item variable vent geometry;
    \item progressive vent opening;
    \item mechanical rupture models;
    \item two-phase and multiphase discharge;
    \item electrolyte flashing;
    \item vent-jet momentum models;
    \item combustion and fire models downstream of the vent.
\end{itemize}